\subsection{Symmetries in the orthogonal group.}

Unless expressly stated otherwise, one assumes, in this no., that A is a commutative field of characteristic ≠ 2, and that Φ is the symmetric bilinear form associated with a nondegenerate quadratic form Q on E. Recall that one has Φ(x, x) = 2Q(x) for x ∈ E (§ 3, no. 4).

Let H be a non-isotropic hyperplane in E, and u the symmetry with respect to H (no. 3). Let \(a \neq 0\) be a vector orthogonal to H; one has by hypothesis \(u(a) = -a\) . Every vector \(x \in E\) can be written in one and only one way in the form \(x = \lambda a + y\) with \(\lambda \in A\) and \(y \in H\) ; since \(a\) and \(y\) are orthogonal, one has \(\Phi(x, a) = \lambda \Phi(a, a)\) , whence, since \(a\) is non-isotropic ( \(\S 4\) , no. 1, cor. of prop. 1), \(\lambda = \Phi(x, a) \Phi(a, a)^{-1}\) . This being so, one has

\[
u (x) = \lambda u (a) + u (y) = - \lambda a + y = x - 2 \lambda a,
\]

whence

\[
u (x) = x - 2 \Phi (x, a) \Phi (a, a) ^ {- 1}. a = x - \Phi (x, a) Q (a) ^ {- 1}. a.\tag{6}
\]

One notes that the last member of (6) retains a meaning when A is a field of characteristic 2, and a is a non-singular vector of E; one verifies immediately that one still has then Q(u(x)) = Q(x) for all x ∈ E, in other words u ∈ O(Q). One also says that the involution u thus defined is the symmetry with respect to the hyperplane orthogonal to a (cf. exerc. 28).

PROPOSITION 5. --- Suppose the vector space E is of finite dimension n. The orthogonal group \(\mathbf{O}(\mathrm{Q})\) associated with Q is then generated by the symmetries with respect to the non-isotropic hyperplanes of E.

The proposition being evident for \(n=0\) , we shall reason by induction on \(n\) . Let \(u\) be an orthogonal transformation of E, and let \(x\) be a non-isotropic vector of E (Lemma 1); we distinguish three cases:

\begin{enumerate}
\def\labelenumi{\alph{enumi})}
\item
  Suppose first that \(u(x) = x\) . Then the hyperplane H orthogonal to \(x\) is non-isotropic, and we have \(u(\mathrm{H}) = \mathrm{H}\) . The restriction \(u'\) of \(u\) to H therefore belongs to the orthogonal group \(\mathbf{O}(Q')\) associated with the restriction \(Q'\) of Q to H. The induction hypothesis implies, since \(Q'\) is nondegenerate, that one has \(u' = v_1' \ldots v_m'\) , where \(v_i'\) is a symmetry with respect to a hyperplane \(L_i\) of H. The endomorphism \(v_i\) of E which extends \(v_i'\) and is such that \(v_i(x) = x\) is then the symmetry with respect to the hyperplane \(Ax + L_i\) of E. We obviously have \(u = v_1v_2 \ldots v_m\) .
\item
  Suppose secondly that \(u(x) = -x\) . If we denote by \(s\) the symmetry with respect to the hyperplane H orthogonal to \(x\) , and if we set \(\nu = su\) , one has \(\nu(x) = x\) , and we are reduced to the case \(a\) ).
\item
  Let us finally pass to the general case, and set \(y = u(x)\) so that \(Q(y) = Q(x)\) . Under these conditions, the vectors \(x - y\) and \(x + y\) cannot both be isotropic, for from the relations \(Q(x - y) = 0\) and \(Q(x + y) = 0\) , one would deduce, by adding memberwise, \(2(Q(x) + Q(y)) = 0\) (§ 3, no. 4, def. 2), whence \(4Q(x) = 0\) , contrary to the hypothesis. Suppose, for example, that \(a = x - y\) is not isotropic; we then have
\end{enumerate}

\[
\Phi (y, a) = \mathrm{Q} (y + a) - \mathrm{Q} (y) - \mathrm{Q} (a) = \mathrm{Q} (x) - \mathrm{Q} (y) - \mathrm{Q} (a) = - \mathrm{Q} (a);
\]

consequently, if one denotes by \(s\) the symmetry with respect to the hyperplane orthogonal to \(a\) , formula (6) proves that \(s(y) = y + a = x\) ; by setting \(\nu = su\) , one has \(\nu(x) = x\) , and we are reduced to the case \(a\) . If \(a = x - y\) is isotropic and \(b = x + y\) not isotropic; one sees in the same way that we are reduced to the case \(b\) .

